\section {Random element of a list (Algorithms)}

This riddle is from Yaron Dorman.

You are given a pointer to a linked list whose length you do not know.
Select one of its elements at random, such that every element has the same chance of being selected,
using $O(1)$ memory and a single pass over the list.

\subsection{Solution}

Walk over the list and keep one selected element. Select the first element.
When you reach the $k$'th element, replace the selected element with it with probability $\frac{1}{k}$.
When the list ends, return the selected element.

Why is this uniform? The $k$'th element is selected when we reach it with probability $\frac{1}{k}$,
and it then has to survive all the following elements: element $j$ does not replace it with probability $1-\frac{1}{j} = \frac{j-1}{j}$.
So for a list of $N$ elements the chance that the $k$'th element is returned is
\begin{center}
$\frac{1}{k} \cdot \frac{k}{k+1} \cdot \frac{k+1}{k+2} \cdots \frac{N-1}{N} = \frac{1}{N}$
\end{center}
which is the same for all the elements.

This is known as reservoir sampling. It also works for selecting $m$ elements: keep the first $m$, and replace a random one of them with the $k$'th element
with probability $\frac{m}{k}$.
